\section{Ripasso e Approfondimento Liste (Lab 14)}

In questo laboratorio rivediamo le operazioni fondamentali sulle liste concatenate e affrontiamo esercizi ricorsivi più complessi.

\subsection{Operazioni di Base (Ripasso)}

\subsubsection{1. Ricerca (listFind)}
\begin{concept}{Trova Primo}
Trova il primo nodo che contiene \texttt{value}.
\end{concept}
\begin{jsbox}
function listFind(head, value) {
    if (!head) return null; // Non trovato
    if (head.val === value) return head; // Trovato!
    return listFind(head.next, value); // Cerca nel resto
}
\end{jsbox}

\subsubsection{2. Inserimento Centrale (listInsert)}
\begin{concept}{Chirurgia dei Puntatori}
Inserisce un nuovo valore SUBITO DOPO il nodo \texttt{x}, aggiornando i puntatori.
\end{concept}
\begin{jsbox}
function listInsert(x, value) {
    if (!x) return;
    
    // 1. Creo il nuovo nodo che punta al vecchio 'next' di x
    let nuovo = { val: value, next: x.next };
    
    // 2. Aggiorno x per puntare al nuovo
    x.next = nuovo;
}
\end{jsbox}

\subsubsection{3. Aggiunta in Coda (listPush con Insert)}
\begin{concept}{Discesa Ricorsiva}
Scende fino all'ultimo nodo e usa \texttt{listInsert}.
\end{concept}
\begin{jsbox}
function listPush(head, value) {
    if (!head) return { val: value, next: null }; // Lista vuota
    
    // Se è l'ultimo nodo (next è null)
    if (!head.next) {
        listInsert(head, value); // Appendo qui
        return head;
    }
    
    listPush(head.next, value); // Scendo ancora
    return head;
}
\end{jsbox}

\subsection{Esercizi Ricorsivi Avanzati}

\subsubsection{1. Verifica Uguaglianza (listEquals)}
\begin{concept}{Confronto Strutturale}
Due liste sono uguali se hanno la stessa lunghezza e gli stessi valori nello stesso ordine.
\end{concept}
\begin{jsbox}
function listEquals(l1, l2) {
    // Caso Base: Entrambe finite contemporaneamente -> OK
    if (!l1 && !l2) return true;
    
    // Caso Base: Una finita e l'altra no -> DIVERSE
    if (!l1 || !l2) return false;
    
    // Controllo valore corrente E ricorsione
    return (l1.val === l2.val) && listEquals(l1.next, l2.next);
}
\end{jsbox}

\subsubsection{2. Rimozione con Taglio (listSplice)}
\begin{concept}{Taglio e Ritorno}
Rimuove \texttt{deleteCount} nodi a partire da \texttt{start}.
\end{concept}
\begin{jsbox}
function listSplice(head, start, deleteCount) {
    if (!head) return null;

    // Caso 1: Non siamo ancora arrivati al punto di taglio
    if (start > 0) {
        // Ricostruisco il nodo corrente e vado avanti
        return {
            val: head.val,
            next: listSplice(head.next, start - 1, deleteCount)
        };
    }

    // Caso 2: Siamo nel punto di taglio (start == 0)
    if (deleteCount > 0) {
        // SALTO questo nodo (non lo ritorno).
        // Chiedo alla ricorsione di rimuoverne ancora (count - 1)
        return listSplice(head.next, 0, deleteCount - 1);
    }

    // Caso 3: Finito di tagliare, ritorno il resto
    return head;
}
\end{jsbox}

\subsubsection{3. Mappa Ricorsiva (listMap)}
\begin{concept}{Trasformazione}
Crea una copia trasformata della lista.
\end{concept}
\begin{jsbox}
function listMap(head, f) {
    if (!head) return null;
    return {
        val: f(head.val),             // Trasformo valore
        next: listMap(head.next, f)   // Ricorro
    };
}
\end{jsbox}
